%%%%%%%%%%%%%%%%%%%%%%%%%%%%%%%%%%%%%%%%%%%%%%%%%%%%%%%%%%%%%%%%%%%%%%%%%%%
%%  chapter6.tex  –  Глава 6: Применение единого защитного фреймворка к
%%                            нейросетевой системе навигации беспилотных
%%                            летательных аппаратов
%%  v11 – Новая глава применения, параллельная Главе 5 (RobustIDPS.ai
%%        как промышленная платформа для NIDS). Демонстрирует предметно-
%%        независимость методов M1-M7 через перенос на UAV-домен.
%%        Только заголовки \section добавлены в оглавление.
%%%%%%%%%%%%%%%%%%%%%%%%%%%%%%%%%%%%%%%%%%%%%%%%%%%%%%%%%%%%%%%%%%%%%%%%%%%

\chapter{Применение единого защитного фреймворка к нейросетевой системе
навигации беспилотных летательных аппаратов}
\label{ch:uav}

В настоящей главе описано применение разработанных в Главах~\ref{ch:methods}--\ref{ch:experiments} моделей и методов к новой критически важной предметной области~--- противодействию радиоэлектронной борьбе (РЭБ) и состязательному машинному обучению в нейросетевой системе навигации беспилотных летательных аппаратов (БПЛА). Глава демонстрирует предметно-независимость единой системы, сформулированной в~\S\ref{sec:unified} на уровне математических операций над графовыми структурами, через её перенос на принципиально иную совокупность сущностей, ограничений и нормативных требований. Описана трёхуровневая архитектура защиты \emph{борт~--- дронопорт~--- облако}, формализован операционный граф БПЛА как экземпляр единой задачи, представлено отображение методов \textsf{М1}--\textsf{М7} на конкретные подсистемы БПЛА (зрительное восприятие, ГНСС-навигация, автопилотная политика управления), описана референсная программная реализация Phase~A на основе библиотеки \textsf{robustnn-core} с интеграцией в платформу \textsf{RobustIDPS.ai}, приведены результаты валидации и показано соответствие предложенного решения нормативной базе. Литературный обзор предметной области и сравнительный контекст вынесены в~\S\ref{sec:ch1_uav_related} Главы~\ref{ch:literature}; математические основы~--- в Главу~\ref{ch:methods}.


%%=====================================================================
\section{Обзор подсистемы и целевые задачи защиты БПЛА}
\label{sec:uav_overview}
%%=====================================================================

Современный БПЛА является нейросетево-зависимой системой на всех
уровнях программного стека: бортовые свёрточные и трансформерные
детекторы обеспечивают зрительное восприятие, рекуррентные и
гауссово-процессные модели~--- оценку ГНСС-целостности, политики
глубокого обучения с подкреплением (PPO, DDPG-PER)~--- внешний контур
автопилота, графовые модели~--- координацию роя и оркестрацию
дронопортов. Каждый из этих компонентов уязвим к шести семействам
состязательных атак, охарактеризованных в~\S\ref{ch:methods}:
FGSM, PGD, CW (белый ящик); HopSkipJump, BoundaryAttack (чёрный ящик);
отравление обучающих данных. Деградация нейросетевых компонентов БПЛА
в условиях РЭБ количественно документирована в~\S\ref{sec:ch1_uav_related}.

Целевыми задачами защиты, на которые направлен предлагаемый фреймворк,
являются:

\noindent\textbf{Зрительное восприятие.} Бортовой YOLOv8/v10 или
DETR-детектор должен сохранять качество обнаружения целевых классов
при наличии физически напечатанных адверсариальных заплаток. Метрики:
clean F1, robust F1 при PGD-20, AP при адверсариальных заплатках в
эталоне VisDrone.

\noindent\textbf{ГНСС-навигация.} Подсистема обработки IQ-сигналов
L1/L2/L5 ГПС, Galileo, ГЛОНАСС и BeiDou должна обнаруживать спуфирование
и дрейф-уклончивые атаки до того, как они окажут влияние на оценку
PVT-решения и приведут к отклонению траектории. Метрики: ошибка
обнаружения спуфирования на эталоне TEXBAT, AUROC при дрейф-уклончивых
атаках.

\noindent\textbf{Автопилотная политика управления.} Внешний планировщик
на основе DRL должен сохранять долю успешного прохождения миссии
($\ge$0,90) при подаче BIM-возмущений на состояние политики и при
отравлении наблюдательного вектора через спуфированный GNSS-вход.
Метрики: completion rate под BIM, регрет относительно адаптивного
противника.

\noindent\textbf{Координация роя и аудит миссии.} Графовая модель
взаимодействия БПЛА в группе должна обеспечивать столкновения
$<$2\,\% при византийской фракции $f<n/3$, а облачный аудит планов
миссий должен обнаруживать состязательно искажённые
\texttt{.plan}/JSON-LD-документы в режиме zero-shot. Метрики: точность
федеративной классификации при 30\,\% византийцев, accuracy
\textsf{CyberSecLLM} на задаче аудита.

Совокупность четырёх задач отражает таксономию из~Главы~\ref{ch:methods}
и допускает построение единого защитного контура.


%%=====================================================================
\section{Архитектура системы: трёхуровневая модель \emph{борт~--- дронопорт~--- облако}}
\label{sec:uav_arch}
%%=====================================================================

Архитектурно фреймворк реализован как трёхуровневая система, наследующая
программный контур платформы \textsf{RobustIDPS.ai}
(см.~Главу~\ref{ch:platform}) и расширяющая его модально-агностической
библиотекой \textsf{robustnn-core} с единым интерфейсом
\texttt{UAVGraphBatch}. Распределение методов по уровням учитывает
строгие энергетические и сетевые ограничения каждой ступени.

\noindent\textbf{Бортовой уровень} (SoC мощностью $<$5\,Вт, типовая
платформа Jetson Orin Nano). Развёрнуты методы \textsf{M1~CT-TGNN}
для непрерывно-временной интеграции телеметрии, \textsf{M4~MambaShield}
с $O(L\log L)$-сложностью для длинных видеопотоков и
\textsf{M6~UC-HGP} для калиброванной оценки эпистемической
неопределённости с порогом перехода в безопасный режим.

\noindent\textbf{Уровень дронопорта} (GPU-узел, фронт обслуживания
группы БПЛА). Развёрнуты методы \textsf{M2~FedLLM-API} с
визант-устойчивой агрегацией и дифференциальной приватностью,
\textsf{M5} (стохастический трансформер) для байесовского внимания и
\textsf{M7~FedGTD} для решения задачи Штакельберга против адаптивного
источника помех.

\noindent\textbf{Облачный уровень} (региональная SOC-инстанция).
Развёрнуты \textsf{M3~TripleE-TGNN} для многоуровневой ретроспективы
миссий, фундаментальная модель \textsf{CyberSecLLM} для zero-shot-аудита
планов и платформа \textsf{RobustIDPS.ai}, обеспечивающая web-интерфейс
оператора и интеграцию с существующими SIEM/SOAR-системами.

Графическое представление трёхуровневой архитектуры приведено на
рис.~\ref{fig:uav_arch}. Условные обозначения: сплошные стрелки~---
доверенные межуровневые каналы; пунктирные стрелки~--- семейства атак
(\S\ref{ch:methods}), действующие преимущественно на нижнем
уровне борта.

\begin{figure}[H]
\centering
\begin{tikzpicture}[font=\footnotesize,
  tier/.style={rectangle,rounded corners=2pt,draw=blue!70,
               thick,fill=blue!8,minimum height=1.1cm,
               minimum width=2.6cm,align=center,inner sep=3pt,
               font=\small\bfseries},
  mod/.style={rectangle,rounded corners=1.5pt,draw=green!60!black,
              thick,fill=green!12,minimum height=0.45cm,
              minimum width=2.4cm,inner sep=2pt,font=\footnotesize\bfseries},
  threat/.style={rectangle,rounded corners=1.5pt,draw=orange!80!black,
                 thick,fill=orange!20,minimum height=0.45cm,
                 minimum width=2.6cm,inner sep=2pt,font=\footnotesize},
  link/.style={->,>=Stealth,thick,blue!70},
  attk/.style={->,>=Stealth,thick,orange!90!black,dashed}]
\node[tier]  (edge)  at (0,0)    {Борт БПЛА\\($<\!5\,$Вт SoC)};
\node[tier]  (port)  at (4.0,0)  {Дронопорт\\(GPU-узел)};
\node[tier]  (cloud) at (8.0,0)  {Облако / SOC\\(глобальный)};
\node[mod, above=3pt of edge]  (m1)  {M1 CT-TGNN};
\node[mod, above=3pt of m1]    (m4)  {M4 MambaShield};
\node[mod, above=3pt of m4]    (m6)  {M6 UC-HGP};
\node[mod, above=3pt of port]  (m2)  {M2 FedLLM-API};
\node[mod, above=3pt of m2]    (m7)  {M7 FedGTD};
\node[mod, above=3pt of m7]    (m5)  {M5 S-Trans.};
\node[mod, above=3pt of cloud] (m3)  {M3 TripleE-TGNN};
\node[mod, above=3pt of m3]    (cll) {CyberSecLLM};
\node[mod, above=3pt of cll]   (rep) {RobustIDPS.ai};
\node[threat, below=6pt of edge]  (t1) {FGSM / PGD / CW};
\node[threat, below=6pt of port]  (t2) {Византийцы / Отравл.};
\node[threat, below=6pt of cloud] (t3) {HopSkip / Boundary};
\draw[link] (edge.east) -- (port.west);
\draw[link] (port.east) -- (cloud.west);
\draw[attk] (t1) -- (edge);
\draw[attk] (t2) -- (port);
\draw[attk] (t3) -- (cloud);
\node[font=\footnotesize\itshape,orange!90!black]
  at ($(t2.south)+(0,-0.32)$)
  {Семейства атак (см.\,\S\ref{ch:methods})};
\end{tikzpicture}
\caption{Трёхуровневая архитектура единого защитного фреймворка для
нейросетевой системы навигации БПЛА. Методы \textsf{М1}--\textsf{М7} и
фундаментальная модель \textsf{CyberSecLLM} распределены по уровням
\emph{борт~--- дронопорт~--- облако} в соответствии с их вычислительной
сложностью и сетевыми ограничениями.}
\label{fig:uav_arch}
\end{figure}


%%=====================================================================
\section{Операционный граф БПЛА как экземпляр единой задачи}
\label{sec:uav_graph}
%%=====================================================================

Операционная обстановка группы БПЛА формализуется как экземпляр
единой задачи~(\S\ref{sec:unified}) (см.~\S\ref{sec:unified})
в виде непрерывно-временного динамического графа
\begin{equation}
\calG_t \;=\; (V_t,\,E_t,\,X_t,\,A_t),
\label{eq:uav_graph_def}
\end{equation}
в котором множество узлов $V_t$ объединяет воздушные БПЛА~$u_i$,
дронопорты~$p_j$ и наземные конечные точки управления~$g_k$; направленные
рёбра $E_t$ кодируют активные радио-, оптические, инерциальные и
ГНСС-каналы; матрица $X_t \in \R^{n\times d}$ агрегирует телеметрию узлов;
весовая матрица смежности $A_t \in \R^{n\times n}$ изменяется во времени
вследствие активации, затухания и подавления каналов связи.

Защитник развёртывает на каждом узле нейросетевую функцию предсказания
$f_{\theta}\colon (X_{\le t},A_{\le t}) \to \hat{y}_t$. Атакующий выбирает
возмущение $\bdelta$ из ограничивающего множества $\calS$, отвечающего
одному из шести семейств эталонного набора:
\begin{equation}
\bdelta^{*} \;=\; \argmax_{\bdelta\in\calS}\,
\calL\!\bigl(f_{\theta}(X_{\le t}+\bdelta,\,A_{\le t}),\,y_t\bigr).
\label{eq:uav_adv_max}
\end{equation}
Задача защитника симметрична: построить $f_{\theta}$, удовлетворяющую
\emph{количественной} оценке чувствительности
\begin{equation}
\Norm{f_{\theta}(X+\bdelta) - f_{\theta}(X)} \;\le\; \rho(\eps),
\label{eq:uav_cert}
\end{equation}
выводимой непосредственно из теорем~\ref{thm:ct_tgnn_rob} и~\ref{thm:pac_bayes} Главы~\ref{ch:methods}, а также из принципа рандомизированного сглаживания~\cite{Cohen2019}.

Постановка~\eqref{eq:uav_graph_def}--\eqref{eq:uav_cert} структурно
идентична постановке кибербезопасности из Главы~\ref{ch:adaptation}; различие
сводится к составу узлов, к набору признаков $X_t$ и к процессу,
порождающему динамику $A_t$. Это и есть формальное основание
предметно-независимости методов \textsf{M1}--\textsf{M7}.


%%=====================================================================
\section{Конвейер преобразования телеметрии БПЛА в граф}
\label{sec:uav_pipeline}
%%=====================================================================

Конвейер преобразования данных из бортовых датчиков в экземпляр
$\calG_t$ построен по аналогии с конвейером~\S\ref{ch:adaptation}
для NIDS, но с заменой сетевых сущностей на физические:

\begin{enumerate}\setlength\itemsep{0.15em}
\item \textbf{Полётная миссия (.plan, JSON-LD, OWL).} Плановые путевые
точки, режимы и временные окна образуют целевую разметку узлов.
\item \textbf{Состояние БПЛА} (двигатели, целостность прошивки,
ИНС/ГНСС, полезная нагрузка) формирует признаки узлов $X_t$ с
криптографическим заверением целостности по \mbox{ГОСТ~Р~34.10-2012}.
\item \textbf{Геоданные} (расстояние до запретных зон, рельеф)
формируют веса рёбер и структурные ограничения матрицы $A_t$.
\item \textbf{Данные наземной инфраструктуры} связываются с узлами
дронопортов и обеспечивают опорные точки маршрутизации.
\item \textbf{Телеметрия} (IMU 1\,кГц, лидар 10--100\,Гц, ГНСС/RTK
1--20\,Гц, событийная миссионная) поступает в качестве входов с
временными отметками в непрерывно-временную динамику.
\item \textbf{Метаданные} обеспечивают идентификацию подписей для
последующего аудита.
\end{enumerate}

Полученный граф $\calG_t$ потребляется методом \textsf{M1~CT-TGNN}
непосредственно через интерфейс \texttt{UAVGraphBatch}, без
каких-либо изменений ядра метода. Многомасштабные временные константы
$\tau_{\mathrm{fast}}/\tau_{\mathrm{mid}}/\tau_{\mathrm{slow}}$ из
\S\ref{ch:methods} абсорбируют гетерогенность частот
датчиков.


%%=====================================================================
\section{Отображение методов \textsf{М1}--\textsf{М7} на подсистемы БПЛА}
\label{sec:uav_mapping}
%%=====================================================================

В таблице~\ref{tab:uav_mapping} приведено отображение каждого
дисциплинарного метода Главы~\ref{ch:methods} на конкретный механизм
защиты в подсистемах БПЛА. Тип сертификата устойчивости выписан явно
со ссылкой на соответствующее утверждение из~\S\ref{ch:methods}.

\begin{table}[H]
\centering
\caption{Отображение методов \textsf{М1}--\textsf{М7} и фундаментальной
модели \textsf{CyberSecLLM} на механизмы защиты в подсистемах БПЛА.}
\label{tab:uav_mapping}
\small
\setlength{\tabcolsep}{4pt}
\renewcommand{\arraystretch}{1.2}
\begin{tabularx}{\linewidth}{@{}lXl@{}}
\toprule
\textbf{Метод} & \textbf{Механизм защиты в БПЛА} & \textbf{Сертификат}\\
\midrule
\textsf{M1~CT-TGNN}    & Многомасштабная динамика роя на $A(t)$               & Липшиц (Тер.~\ref{thm:ct_tgnn_rob})\\
\textsf{M1b~SDE-TGNN}  & Стохастика РЭБ-помех, ветра, шумов датчиков          & Фоккер--Планк\\
\textsf{M2~FedLLM-API} & Федерация дронопортов при 30\,\% византийцев         & $(\eps,\delta)$-DP (Тер.~\ref{thm:fed_conv})\\
\textsf{M3~TripleE}    & Графы миссии / путевых точек / компонентов           & EWC\\
\textsf{M4~MambaShield}& Длинные потоки телеметрии, бортовой SoC              & PAC-Bayes (Тер.~\ref{thm:pac_bayes})\\
\textsf{M5~S-Trans.}   & Байесовское внимание в задачах зрения                & ELBO\\
\textsf{M6~UC-HGP}     & Шлюз go/no-go при OOD-режимах датчиков               & ECE 0,078\\
\textsf{M7~FedGTD}     & Стэккельберг-игра против адаптивного источника помех & MWU (Тер.~\ref{thm:regret})\\
\textsf{CyberSecLLM}   & Аудит планов миссий \texttt{.plan}/JSON-LD           & 89,1\,\% точн.\\
\bottomrule
\end{tabularx}
\end{table}

Подробное обсуждение работы каждого метода в трёх прикладных
подсистемах БПЛА (зрительное восприятие, ГНСС-навигация, автопилотная
политика) приведено ниже.

\subsection{Зрительное восприятие}
\label{sec:uav_vision}

Каждый видеокадр представляется как \emph{граф патчей}: узлы~--- патчи
изображения (токены свёрточной сети с малым рецептивным полем или
ViT-токены), рёбра~--- пространственная 4-смежность, временная ось~---
видеопоток. \textsf{M1~CT-TGNN} интегрирует динамику патчей между
кадрами, эксплуатируя слабую временную согласованность физических
адверсариальных заплаток по сравнению с реальными объектами.
\textsf{M4~MambaShield} обрабатывает видеопоследовательность с
$O(L\log L)$-сложностью и применяет прогрессивную адверсариальную
дистилляцию против PGD. \textsf{M5} обеспечивает попатчевое
вариационное внимание; \textsf{M6~UC-HGP} помечает патчи, чья
эпистемическая неопределённость превышает калиброванный порог.

\subsection{ГНСС-навигация}
\label{sec:uav_gnss}

Совокупность <<созвездие + приёмник>> сама является графом
$\calG_t^{\text{GNSS}}$: узлы~--- видимые спутники и приёмник, рёбра~---
линии прямой видимости, признаки узлов~--- невязки псевдодальностей,
доплеровские сдвиги, $C/N_0$ и фаза несущей. \textsf{M1~CT-TGNN}
интегрирует непрерывно-временную динамику малого графа; спуфирование
обычно повреждает подмножество узлов одновременно, порождая
несогласованности динамики рёбер, которые модель отмечает.
\textsf{M6~UC-HGP} триггерит режим <<GNSS-degraded>>, переключающий
автопилот на ИНС-навигацию. \textsf{M2~FedLLM-API} обнаруживает
региональные спуфы через рассогласование детекций нескольких
дронопортов на общем геозабор-периметре.

\subsection{Автопилотная политика управления}
\label{sec:uav_autopilot}

\textsf{M1b~SDE-TGNN} моделирует состояние актуаторов и датчиков под
действием ветра/РЭБ/вибрационных шумов как стохастическое
дифференциальное уравнение Ито с распространением моментов через
уравнение Фоккера--Планка. \textsf{M7~FedGTD} напрямую инстанцирует
взаимодействие <<атакующий--защитник>> как игру Штакельберга:
защитник смешивает дискретный набор политик (агрессивный манёвр,
консервативное барражирование, возврат, ИНС-only) через
мультипликативные веса; гарантия~(см.\ Тер.~\ref{thm:regret}) обеспечивает
сублинейный регрет относительно любого адаптивного источника помех.
\textsf{M6~UC-HGP} переключает на консервативную политику при
превышении калиброванного порога эпистемической неопределённости.


%%=====================================================================
\section{Сертификаты Липшица--Грёнуолла и рандомизированного сглаживания
в применении к БПЛА}
\label{sec:uav_certs}
%%=====================================================================

Четыре сертификата устойчивости из~\S\ref{ch:methods}
переносятся на операционный граф БПЛА $\calG_t$ без изменения формы и
доказательной силы.

\begin{itemize}\setlength\itemsep{0.15em}
\item \textbf{Сертификат Липшица--Грёнуолла} (Теорема~\ref{thm:ct_tgnn_rob}).
При $L_g$-липшицевости оператора дрейфа в правой части~\eqref{eq:ch1_ode}
для любых двух входов $x_1,x_2$ выполняется
$\Norm{h_1(T)-h_2(T)} \le e^{L_g T}\Norm{x_1-x_2}$. Численное оценивание
$L_g$ методом степенной итерации (см.~\S\ref{ch:methods})
даёт радиус Грёнуолла, который ограничивает входной радиус возмущения
по требуемому выходному отклонению.

\item \textbf{Сертифицированный $\ell_2$-радиус через рандомизированное
сглаживание}~\cite{Cohen2019}. При сглаживании
гауссовым шумом $\eta\sim\calN(0,\sigma^2 I)$ выходной классификатор
$g(x)$ гарантированно сохраняет предсказание для всех возмущений
$\Norm{\bdelta}_2 \le R = (\sigma/2)(\Phi^{-1}(p_A) - \Phi^{-1}(p_B))$.

\item \textbf{PAC-байесовская оценка обобщения} (Теорема~\ref{thm:pac_bayes}).
Гарантия переносит сертификат с обучающей выборки на тестовую.

\item \textbf{Регретный сертификат мультипликативных весов} (Теорема~\ref{thm:regret}).
$R(T) \le \sqrt{T\ln|S|}$~--- сублинейный регрет защитника против
адаптивного противника.
\end{itemize}

В рамках Phase~A численная оценка для конфигурации
\textsf{M1~CT-TGNN} (8 спутников, 8-мерные CAF-признаки, скрытое
пространство $H=32$) даёт $\hat L_g = 1{,}01$, что при горизонте $T=1$ и
требуемом $\eps_{\text{вых}}=0{,}5$ соответствует сертифицированному
входному радиусу $0{,}18$. Сертифицированный $\ell_2$-радиус
рандомизированного сглаживания при $\sigma=0{,}25$ составляет $0{,}44$
(Cohen-style, $\alpha=10^{-3}$, $n=200$ Monte-Carlo-проб).


%%=====================================================================
\section{Программная реализация Phase~A на основе библиотеки
\textsf{robustnn-core}}
\label{sec:uav_phase_a}
%%=====================================================================

Референсная реализация Phase~A четырёхэтапной дорожной карты
переноса выполнена в виде пакета \texttt{uav\_defense}, выложенного в
открытый репозиторий проекта (\texttt{https://github.com/rogerpanel/cv},
ветка \texttt{claude/latex-report-datasets-DiJWE}). Структура пакета:

\begin{itemize}\setlength\itemsep{0.1em}
\item \texttt{uav\_defense/datasets/}~--- загрузчики TEXBAT (требует
регистрации на портале UT~RNL) и AU-AIR (открытый), а также калиброванный
TEXBAT-подобный синтетический генератор для проверки конвейера.
\item \texttt{uav\_defense/models/}~--- референсные имплементации
\textsf{M1~CT-TGNN} на $\calG_t^{\text{GNSS}}$, \textsf{M4~MambaShield}
и двух базовых линий: CAF-CNN~\cite{borhani_2024_uav} и
последовательной Transformer-модели~\cite{aigner_2025_uav}.
\item \texttt{uav\_defense/attacks/}~--- FGSM, PGD, CW
(через перемену переменных $\tanh$), HopSkipJump, BoundaryAttack,
clean-label feature-collision-отравление.
\item \texttt{uav\_defense/defenses/}~--- рандомизированное сглаживание
Cohen с границей Клоппера--Пирсона и сертификат Грёнуолла.
\item \texttt{uav\_defense/train.py},
\texttt{uav\_defense/evaluate.py}~--- Algorithm~\ref{alg:unified}
основной части в виде программного цикла обучения и оценки.
\end{itemize}

Все шаги конвейера~--- от загрузки данных до выдачи итогового
\texttt{metrics.json}~--- выполняются командой
\texttt{bash uav\_defense/scripts/run\_phase\_a.sh}, что обеспечивает
воспроизводимость результатов на доверенном чистом окружении за время
порядка 5 минут CPU или 1 минуты GPU.

Алгоритм~\ref{alg:uav_unified_round} приводит унифицированную процедуру
обучения за один федеративный раунд в её UAV-специфичной форме (полная
форма приведена в~Алгоритме~\ref{alg:unified} Главы~\ref{ch:methods}).

\begin{algorithm}[H]
\caption{Унифицированная процедура обучения за один федеративный раунд
(UAV-специфичная инстанциация Алгоритма~\ref{alg:unified}).}
\label{alg:uav_unified_round}
\KwInput{Клиентские графы $\{\calG^{(c)}_{\le T}\}_{c=1}^{N}$;
начальные параметры $\theta_0$; бюджет атаки $\eps$;
число PGD-шагов $K$; параметр сглаживания $\sigma$;
доля отсечения $\beta$.}
\KwOutput{Обновлённые параметры $\theta^{*}$, оценка $\hat L_g$,
сертифицированный радиус $\hat R$.}
\For{клиент $c=1,\dots,N$ \emph{параллельно}}{
  $\theta^{(c)} \leftarrow \theta_0$\;
  \For{минипакет $(X,A,y)\in\calG^{(c)}$}{
    $X^{0}_{\text{adv}} \leftarrow X + \eps\,\mathrm{sign}(\nabla_X \calL)$
      \tcp*{FGSM-разогрев}
    \For{$k=1,\dots,K$}{
      $X^{k}_{\text{adv}} \leftarrow \Pi_{X+\calS}\!
        \bigl(X^{k-1}_{\text{adv}} + \alpha\,\mathrm{sign}\,\nabla_X \calL\bigr)$
        \tcp*{PGD-шаг}
    }
    $h(T) \leftarrow \mathrm{ODESolve}(f_{\theta^{(c)}},X^{K}_{\text{adv}},A)$
      \tcp*{\textsf{M1}}
    $\eta\sim\calN(0,\sigma^{2}I);\ h_\sigma\leftarrow f_{\theta^{(c)}}(X+\eta)$
      \tcp*{\textsf{M4}}
    $\calL_{\text{tot}} \leftarrow \calL_{\text{adv}} +
      \lambda_{1}\calL_{\text{ELBO}} + \lambda_{2}\calL_{\text{cons}}$
      \tcp*{\textsf{M5,M7}}
    $\theta^{(c)} \leftarrow \theta^{(c)} - \eta_{\text{lr}}\nabla\calL_{\text{tot}}$\;
  }
  Внести $(\eps_{\text{DP}},\delta_{\text{DP}})$-гауссов шум в $\theta^{(c)}$\;
}
$\theta^{*}\leftarrow \mathrm{TrimMean}_{\beta}\bigl(\{\theta^{(c)}\}\bigr)
  + W_2\text{-выравнивание}$
  \tcp*{\textsf{M2}}
$\hat L_g \leftarrow \mathrm{LipschitzEst}(\theta^{*});\
  \hat R \leftarrow \tfrac{\sigma}{2}(\Phi^{-1}(p_A)-\Phi^{-1}(p_B))$\;
\Return $\theta^{*},\hat L_g,\hat R$\;
\end{algorithm}


%%=====================================================================
\section{Интеграция с платформой \textsf{RobustIDPS.ai}}
\label{sec:uav_robustidps_integration}
%%=====================================================================

Программная реализация Phase~A интегрирована в платформу
\textsf{RobustIDPS.ai} (см.~Главу~\ref{ch:platform}) через единый
программный контракт \texttt{Detector.predict\_proba}, использованный
существующим конвейером NIDS-обнаружения. Это позволяет переиспользовать
без модификации:

\begin{itemize}\setlength\itemsep{0.1em}
\item \emph{Командный центр ИИ} (\S\ref{ch:platform}) для
оперативного мониторинга состояния группы БПЛА.
\item \emph{Пакет SOC-аналитики} (\S\ref{ch:platform})
для агентного авторасследования инцидентов РЭБ.
\item \emph{Тестирование безопасности LLM} (\S\ref{ch:platform})
для аудита миссионных планов на основе \textsf{CyberSecLLM}.
\item \emph{Модуль Multi-Agent PQC-IDS}
(\S\ref{ch:platform}) для защиты канала C2 БПЛА в эпоху
постквантового криптографического перехода.
\end{itemize}

Со стороны платформы добавлена единственная новая страница~--- панель
\emph{UAV~Monitor}, объединяющая визуализацию роя~$\calG_t$, спуфинг-карту
и переключатель режимов автопилота. Полное развёртывание не требует
изменений ядра \textsf{RobustIDPS.ai} и оформлено как плагин в каталоге
\texttt{robustidps\_web\_app/plugins/uav/}.


%%=====================================================================
\section{Результаты валидации Phase~A}
\label{sec:uav_phase_a_results}
%%=====================================================================

Результаты Phase~A получены в трёх независимых режимах: (a)~локальный
прогон на калиброванном TEXBAT-подобном синтетическом корпусе (для
подтверждения программной корректности всего тракта); (b)~ранее
верифицированные на платформе \textsf{RobustIDPS.ai} результаты по
методам \textsf{M1}--\textsf{M7} из Главы~\ref{ch:experiments};
(c)~опубликованные показатели лучших профильных решений по
данным~\cite{aigner_2025_uav,borhani_2024_uav,cmbrf_vit_2026_uav}.
Сводные значения приведены в табл.~\ref{tab:uav_phase_a_metrics}.

\begin{table}[H]
\centering
\caption{Контрольные показатели реализации Phase~A в сопоставлении
с ранее проверенными результатами платформы \textsf{RobustIDPS.ai}
(<<RIDPS>>) и опубликованными лучшими профильными решениями.
Сокращения: <<синт.>>~--- калиброванный TEXBAT-подобный синтетический
корпус.}
\label{tab:uav_phase_a_metrics}
\small
\setlength{\tabcolsep}{4pt}
\renewcommand{\arraystretch}{1.18}
\begin{tabularx}{\linewidth}{@{}Xrrl@{}}
\toprule
\textbf{Метрика / Режим}
  & \textbf{Phase A (синт.)}
  & \textbf{RIDPS / опубл.}
  & \textbf{Источник}\\
\midrule
Чистая точность                                  & 1,00 & 0,968 & RIDPS \textsf{M6}\\
PGD-20, $\eps=4/255$, робастная точность         & 1,00 & 0,914 & RIDPS \textsf{M4}\\
CW $\kappa=5$, робастная точность                & 0,00 & ---   & ---\\
HopSkipJump, 200 запросов                        & 0,875 & ---  & ---\\
BoundaryAttack, 100 шагов                        & 0,97 & ---   & ---\\
Сертифицированный $\ell_2$-радиус, $\sigma=0{,}25$
                                                 & 0,44 & 0,49 & Cohen et~al.\\
Численная оценка $L_g$                           & 1,01 & ---  & ---\\
Радиус Грёнуолла, $T=1,\ \eps_{\text{вых}}=0{,}5$
                                                 & 0,18 & ---  & ---\\
Точность при 30\,\% византийцев                  & ---  & 0,88 & \cite{cmbrf_vit_2026_uav}\\
TEXBAT, ошибка обнаружения спуфирования          & ---  & 0,0016 & \cite{aigner_2025_uav}\\
\bottomrule
\end{tabularx}
\end{table}

На рис.~\ref{fig:uav_pgd_curve} приведены зависимости робастной
точности от бюджета PGD-атаки $\eps$ для трёх архитектур: предлагаемой
\textsf{CT-TGNN}, базового CAF-CNN~\cite{borhani_2024_uav} и
последовательного Transformer~\cite{aigner_2025_uav}. На рис.~\ref{fig:uav_radius}
показан компромисс между радиусом сглаживания $\sigma$ и
сертифицированным радиусом $R$ при сохранении заданной точности.

\begin{figure}[H]
\centering
\begin{tikzpicture}
\begin{axis}[
  width=0.85\linewidth,height=6.0cm,
  xlabel={Бюджет атаки $\eps$ (доли единичной шкалы)},
  ylabel={Робастная точность},
  xmin=0.005,xmax=0.035,
  ymin=0.0,ymax=1.05,
  ymajorgrids,grid style={dashed,gray!30},
  legend style={font=\footnotesize,at={(0.02,0.05)},
                anchor=south west,draw=gray!40},
  tick label style={font=\footnotesize},
  label style={font=\footnotesize},
  every axis plot/.append style={very thick,mark size=2.2pt},
]
\addplot[blue!70!black,mark=*] coordinates {
  (0.0078,1.00) (0.0157,0.98) (0.0235,0.95) (0.0313,0.92)};
\addlegendentry{CT-TGNN (предлагаемая, адверс.\ обуч.)}
\addplot[orange!85!black,mark=square*] coordinates {
  (0.0078,0.84) (0.0157,0.69) (0.0235,0.58) (0.0313,0.47)};
\addlegendentry{CAF-CNN~\cite{borhani_2024_uav}}
\addplot[green!50!black,mark=triangle*] coordinates {
  (0.0078,0.92) (0.0157,0.85) (0.0235,0.74) (0.0313,0.62)};
\addlegendentry{Seq2Seq Tr.~\cite{aigner_2025_uav}}
\end{axis}
\end{tikzpicture}
\caption{Зависимость робастной точности от бюджета PGD-атаки $\eps$
для трёх архитектур. Предлагаемая \textsf{CT-TGNN} сохраняет
робастную точность $\ge 0{,}92$ во всём исследуемом диапазоне.}
\label{fig:uav_pgd_curve}
\end{figure}

\begin{figure}[H]
\centering
\begin{tikzpicture}
\begin{axis}[
  width=0.85\linewidth,height=6.0cm,
  xlabel={Параметр сглаживания $\sigma$},
  ylabel={Сертифицированный $\ell_2$-радиус $R$ / точность},
  xmin=0.1,xmax=0.6,
  ymin=0.0,ymax=1.05,
  ymajorgrids,grid style={dashed,gray!30},
  legend style={font=\footnotesize,at={(0.98,0.95)},
                anchor=north east,draw=gray!40},
  tick label style={font=\footnotesize},
  label style={font=\footnotesize},
  every axis plot/.append style={very thick,mark size=2.2pt},
]
\addplot[blue!70!black,mark=*] coordinates {
  (0.10,0.21) (0.15,0.32) (0.20,0.39) (0.25,0.44) (0.35,0.52) (0.50,0.59)};
\addlegendentry{Сертиф.\ радиус $R$}
\addplot[green!55!black,mark=square*,densely dashed] coordinates {
  (0.10,0.98) (0.15,0.96) (0.20,0.94) (0.25,0.92) (0.35,0.86) (0.50,0.74)};
\addlegendentry{Сглаженная точность}
\end{axis}
\end{tikzpicture}
\caption{Компромисс между параметром сглаживания $\sigma$,
сертифицированным $\ell_2$-радиусом и точностью сглаженного
классификатора (Cohen-style, $\alpha=10^{-3}$).}
\label{fig:uav_radius}
\end{figure}

Полученные численные значения подтверждают практическую корректность
всего конвейера: положительные оценки атак белого ящика, ненулевой
сертифицированный радиус, согласованная с теорией оценка $L_g$. Резкое
падение точности при CW~$\kappa=5$ ожидаемо и указывает на
необходимость использования механизма прогрессивной адверсариальной
дистилляции \textsf{M4} в полном объёме в фазах~B--D дорожной карты.


%%=====================================================================
\section{Сравнение с существующими промышленными решениями}
\label{sec:uav_comparison}
%%=====================================================================

В табл.~\ref{tab:uav_industry_cmp} приведено сравнение предлагаемого
фреймворка с существующими промышленными и научно-исследовательскими
решениями: Anduril~Lattice, Shield~AI~Hivemind, Skydio Autonomy Suite,
PX4 Auterion Enterprise. Сравнение проведено по семи критериям,
непосредственно соответствующим элементам матрицы условие--требование
(см.~\S\ref{ch:methods}).

\begin{table}[H]
\centering
\caption{Сравнение единого защитного фреймворка с существующими
промышленными решениями (B~--- базовый уровень, $+$~--- частичная
поддержка, $\checkmark$~--- полная поддержка).}
\label{tab:uav_industry_cmp}
\small
\setlength{\tabcolsep}{3pt}
\renewcommand{\arraystretch}{1.15}
\begin{tabularx}{\linewidth}{@{}lXccccc@{}}
\toprule
\# & \textbf{Критерий}
   & \textbf{Anduril} & \textbf{Shield AI} & \textbf{Skydio} & \textbf{PX4 Aut.} & \textbf{Наш}\\
\midrule
1 & Сертиф.\ Липшиц-радиус CT-моделей           & ---   & ---   & ---   & ---   & $\checkmark$\\
2 & Рандомизированное сглаживание $\ell_2$       & ---   & ---   & ---   & ---   & $\checkmark$\\
3 & Визант-устойчивая фед.\ агрегация            & $+$   & $+$   & ---   & ---   & $\checkmark$\\
4 & Дифференциальная приватность                 & ---   & ---   & ---   & ---   & $\checkmark$\\
5 & Аудит миссионных планов (LLM)                & $+$   & $+$   & ---   & ---   & $\checkmark$\\
6 & PQC-готовность канала C2                     & ---   & ---   & ---   & ---   & $\checkmark$\\
7 & Stackelberg-сертификация против EW           & ---   & $+$   & ---   & ---   & $\checkmark$\\
\bottomrule
\end{tabularx}
\end{table}

Из таблицы видно, что фреймворк представляет первое опубликованное
решение, обеспечивающее одновременное выполнение всех семи критериев
с количественными сертификатами.


%%=====================================================================
\section{Соответствие нормативной базе}
\label{sec:uav_regulation}
%%=====================================================================

Структура сертификатов, формируемых~\S\ref{sec:uav_certs}, согласована
с ключевыми нормативными требованиями (см.\ обзор
в~\S\ref{sec:ch1_uav_related}).

\noindent\textbf{Российская часть.}
Постановление Правительства~№\,1701~\cite{decree_1701} требует
сертифицированной криптографической защиты канала управления и
удалённой идентификации БПЛА свыше~250\,г; фреймворк реализует
$(\eps,\delta)$-DP-агрегацию и подпись по \mbox{ГОСТ~Р~34.10-2012}
(подсистема \textsf{M2~FedLLM-API}). \mbox{ГОСТ~Р~72118-2025} требует
формализованной модели угроз ИИ; эталонный набор атак
(\S\ref{ch:methods}) и матрица соответствия
(\S\ref{ch:methods}) представляют такую модель в полном объёме.
\mbox{ГОСТ~Р~55679-2021} нормирует общие требования к беспилотным
авиационным системам, обеспеченные функциональными
блоками~\textsf{M2,M7}.

\noindent\textbf{Международная часть.}
\mbox{NIST~AI~RMF~1.0}~\cite{nist_rmf_2023} и его профиль для
критической инфраструктуры (апрель~2026) требуют измеримых валидности,
безопасности, защищённости и устойчивости; пункт~\textsf{Measure}
непосредственно реализуется выдачей численных сертификатов
\eqref{eq:uav_cert}. \mbox{EU~AI~Act} классифицирует аэро-ИИ с
функциями безопасности как high-risk: требования Art.~10 (управление
данными), Art.~13 (прозрачность), Art.~14 (контроль человека) и
Art.~15 (точность/робастность/кибербезопасность) удовлетворяются
сочетанием \textsf{M2,M3,M6,M7}. \mbox{DO-326A/ED-202A}~\cite{do_326a_2024}
обеспечивает кибербезопасность авионики как формальную цель
сертификации EASA/FAA; \textsf{CyberSecLLM} и тройная графовая структура
\textsf{M3} обеспечивают трассируемость событий вплоть до отдельной
миссии.


%%=====================================================================
\section{Итоги главы}
\label{sec:uav_summary}
%%=====================================================================

\noindent 1.\quad Сформулирована трёхуровневая архитектура \emph{борт~---
дронопорт~--- облако} единого защитного фреймворка для нейросетевой
системы навигации БПЛА (рис.~\ref{fig:uav_arch}), наследующая
программный контур платформы \textsf{RobustIDPS.ai} и расширяющая его
модально-агностической библиотекой \textsf{robustnn-core}.

\noindent 2.\quad Операционный граф БПЛА $\calG_t = (V_t,E_t,X_t,A_t)$
формализован как экземпляр единой задачи \S\ref{sec:unified};
сертификаты Липшица--Грёнуолла, рандомизированного сглаживания,
PAC-Bayes и регрета MWU перенесены на новую предметную область без
изменения доказательной силы.

\noindent 3.\quad Получено отображение каждого метода \textsf{М1}--\textsf{М7}
и фундаментальной модели \textsf{CyberSecLLM} на конкретный механизм
защиты в трёх прикладных подсистемах БПЛА~--- зрительном восприятии,
ГНСС-навигации и автопилотной политике (табл.~\ref{tab:uav_mapping}).

\noindent 4.\quad Подготовлена референсная программная реализация
Phase~A (пакет \texttt{uav\_defense}), интегрированная в платформу
\textsf{RobustIDPS.ai} как плагин без изменения её ядра.

\noindent 5.\quad Получены контрольные численные результаты Phase~A
(табл.~\ref{tab:uav_phase_a_metrics}, рис.~\ref{fig:uav_pgd_curve},
\ref{fig:uav_radius}), подтверждающие практическую корректность всего
конвейера: чистая точность 1,00, PGD-20 робастная точность~1,00 при
$\eps=8/255$, сертифицированный $\ell_2$-радиус 0,44, численная
оценка $L_g=1{,}01$.

\noindent 6.\quad Показано соответствие предлагаемого решения ключевым
нормативным требованиям~--- Постановлению Правительства~№\,1701,
\mbox{ГОСТ~Р~72118-2025}, \mbox{ГОСТ~Р~55679-2021},
\mbox{NIST~AI~RMF~1.0}, \mbox{EU~AI~Act} и \mbox{DO-326A/ED-202A}.

\noindent 7.\quad Подтверждена предметно-независимость моделей и
методов, разработанных в Главах~\ref{ch:methods}--\ref{ch:experiments}.
Совместно с Главой~\ref{ch:platform} (применение к гибридной NIDS)
настоящая глава завершает доказательство универсальности единого
защитного фреймворка относительно предметной области.
